\chapter{The Acme Screen}
\label{ch:screen}

\index{screen layout}

When acme starts, you see a deceptively simple display.  There are no
menus, no toolbars, no status bar, no file tree.  There is only text,
arranged in columns and windows.  Every element on screen is either
text you can edit or a border you can drag.

\section{Anatomy of the Display}

The acme screen is divided into a hierarchy:

\begin{enumerate}
\item A \textbf{top-level tag bar} spans the full width of the screen.
\item Below it, one or more \textbf{columns} are arranged side by side.
\item Each column contains one or more \textbf{windows}, stacked vertically.
\item Each window has its own \textbf{tag bar} and \textbf{body}.
\end{enumerate}

\subsection{The Top Tag Bar}
\index{tag bar!top-level}

The top tag bar shows global information and commands.  A typical top
tag bar reads:

\begin{Verbatim}[frame=single]
Newcol Kill Putall Dump Exit
\end{Verbatim}

These are not buttons---they are editable text.  You can add your own
commands by typing them into the tag bar.  Middle-click (\button{2}) on
any word to execute it.

\subsection{Columns}
\index{columns}

Columns are vertical panes that divide the screen horizontally.  Each
column has a thin tag bar at its top, typically containing:

\begin{Verbatim}[frame=single]
New Cut Paste Snarf Sort Zerox Delcol
\end{Verbatim}

\noindent
You can:

\begin{itemize}
\item Create a new column by executing \cmd{Newcol} in the top tag bar.
\item Delete a column by executing \cmd{Delcol} in the column's tag bar.
\item Resize columns by dragging the thin border between them with
      \button{1}.
\end{itemize}

\subsection{Windows}
\index{windows}

Each window has two parts:

\paragraph{The tag bar}
\index{tag bar!window}
is a single line (or more, if the content wraps) at the top of the
window.  It contains the window's filename followed by commands.  A
typical tag bar reads:

\begin{Verbatim}[frame=single]
/home/user/src/main.c Del Snarf Get Put Undo Redo | Look
\end{Verbatim}

The filename at the left is editable---you can change it to rename the
file, or clear it for a scratch buffer.  Everything after the vertical
bar \texttt{|} is the ``user area'' of the tag, where you can type
whatever commands you find useful.

\paragraph{The body}
\index{body}
occupies the rest of the window and contains the file's text.  This is
where you read and edit.

\subsection{The Layout Button}
\index{layout button}

At the far left of each window's tag bar is a small square.  This is
the \emph{layout button}.  It serves two purposes:

\begin{itemize}
\item Drag it with \button{1} to move the window within its column or
      to another column.
\item Its appearance indicates the window's state:
  \begin{itemize}
  \item A clean square means the window matches the file on disk.
  \item A filled (dirty) square means the window has unsaved changes.
  \end{itemize}
\end{itemize}

\subsection{The Scrollbar}
\index{scrollbar}

Along the left edge of each window body is a narrow scrollbar.  It
shows a small rectangle (the ``thumb'') indicating which portion of the
file is currently visible.  See Section~\ref{sec:scrollbar} for details
on mouse interaction with the scrollbar.

\section{Text Everywhere}

The key insight of acme's interface is that \emph{everything is text}.
The tag bars are text.  The body is text.  Commands are text.  You
interact with all of it the same way: by clicking on it with the mouse.

This means:

\begin{itemize}
\item You can type a command anywhere---in a tag, in a body, even in
      the middle of a file---and middle-click it to execute.
\item You can type a filename in any window and right-click it to open
      that file.
\item Output from commands appears as text in windows, which you can
      then edit, search, or use as input to other commands.
\end{itemize}

There is no distinction between ``code,'' ``output,'' and ``interface.''
It is all text, and it is all interactive.

\section{Window Management}

\index{window management}

Windows in acme are not managed by tabs or a file tree.  Instead:

\begin{itemize}
\item Open a file by right-clicking (\button{3}) its name.
\item Create an empty window with the \cmd{New} command.
\item Close a window with the \cmd{Del} command.
\item Clone a window with the \cmd{Zerox} command.
\item Sort windows alphabetically with the \cmd{Sort} command.
\item Move windows by dragging the layout button.
\end{itemize}

To navigate between windows, you simply click in the one you want.
There is no ``active window'' highlight or focus indicator---you click
where you want to work, and that window responds.

\section{Directory Windows}
\index{directory windows}

Right-clicking a directory path opens a window listing the directory's
contents.  Each entry in the listing is itself clickable: right-click a
filename to open it, right-click a subdirectory to list it.

Directory windows update automatically when files change.  They serve
the role of a file browser, but they are just text in a window---you
can search them, select from them, or pipe their contents through
commands.

\section{The +Errors Window}
\index{+Errors window}
\index{errors}

When an external command writes to standard error, the output appears in
a window named \file{+Errors} (prefixed with the directory path).
These windows are created automatically as needed.  You can close them
with \cmd{Del} like any other window.

Compiler errors often appear in the format \file{file.c:27: error...}.
Right-clicking (\button{3}) on such a line opens \file{file.c} at
line~27---a simple but powerful way to navigate error output.
